% --- Archon physics formalization source begin ---
% archon:physics
% archon:covers PhyXMiniProblems/problem_phyx_mini_0826.lean
% archon:source-report reports/phyx_mini/problem_phyx_mini_0826.source.json

\chapter{Physics problem phyx\_mini\_0826}
\label{ch:PhyXMiniProblems_problem_phyx_mini_0826}

\paragraph{Problem source.}
\#\# Physical scenario

Figure shows two disks: an engine flywheel (A) and a clutch plate (B) attached to a transmission shaft. Their moments of inertia are \textbackslash{}( I\_A \textbackslash{}) and \textbackslash{}( I\_B \textbackslash{}); initially, they are rotating in the same direction with constant angular speeds \textbackslash{}( \textbackslash{}omega\_A \textbackslash{}) and \textbackslash{}( \textbackslash{}omega\_B \textbackslash{}), respectively. We push the disks together with forces acting along the axis, so as not to apply any torque on either disk. The disks rub against each other and eventually reach a common angular speed \textbackslash{}( \textbackslash{}omega \textbackslash{}).

\#\# Question

Derive an expression for \textbackslash{}( \textbackslash{}omega \textbackslash{}).

\#\# Answer choices

- A: A: \textbackslash{}( \textbackslash{}omega = \textbackslash{}frac\{I\_A \textbackslash{}omega\_A + I\_B \textbackslash{}omega\_B\}\{I\_A - I\_B\} \textbackslash{})

- B: B: \textbackslash{}( \textbackslash{}omega = \textbackslash{}frac\{I\_A \textbackslash{}omega\_A + I\_B \textbackslash{}omega\_B\}\{I\_A + I\_B\} \textbackslash{})

- C: C: \textbackslash{}( \textbackslash{}omega = \textbackslash{}frac\{I\_A \textbackslash{}omega\_B + I\_B \textbackslash{}omega\_A\}\{I\_A + I\_B\} \textbackslash{})

- D: D: \textbackslash{}( \textbackslash{}omega = \textbackslash{}frac\{I\_A \textbackslash{}omega\_B + I\_B \textbackslash{}omega\_A\}\{I\_A - I\_B\} \textbackslash{})

\#\# Figure caption (auxiliary; use the image as primary evidence)

The image consists of two diagrams labeled "BEFORE" and "AFTER", depicting two mechanical systems.

**BEFORE:**

- On the left: A large gear labeled \textbackslash{}( I\_A \textbackslash{}) with an arrow indicating counterclockwise angular velocity \textbackslash{}( \textbackslash{}omega\_A \textbackslash{}). A force \textbackslash{}( \textbackslash{}vec\{F\} \textbackslash{}) is applied to its left side.
- On the right: A smaller disk labeled \textbackslash{}( I\_B \textbackslash{}) with an arrow showing clockwise angular velocity \textbackslash{}( \textbackslash{}omega\_B \textbackslash{}). A force \textbackslash{}( -\textbackslash{}vec\{F\} \textbackslash{}) is applied to its right side.
- The gear and disk are connected along a dashed line indicating their axes are aligned.

**AFTER:**

- The gear and disk are combined into a single unit labeled \textbackslash{}( I\_A + I\_B \textbackslash{}).
- An arrow indicates a counterclockwise angular velocity \textbackslash{}( \textbackslash{}omega \textbackslash{}).
- The same forces \textbackslash{}( \textbackslash{}vec\{F\} \textbackslash{}) and \textbackslash{}( -\textbackslash{}vec\{F\} \textbackslash{}) are applied to the combined object, similar to the "BEFORE" scenario.

The diagrams represent a change from separate rotational bodies to a combined rotational system under applied forces.

\#\# Dataset metadata

Rotational Motion; Physical Model Grounding Reasoning, Multi-Formula Reasoning

\paragraph{Recorded answer/context.}
B: B: \textbackslash{}( \textbackslash{}omega = \textbackslash{}frac\{I\_A \textbackslash{}omega\_A + I\_B \textbackslash{}omega\_B\}\{I\_A + I\_B\} \textbackslash{})

\paragraph{Figure/image path.}
/root/proposal\_for\_physic/science-mango/phyx\_mini\_run/phyx\_data/test\_image/826.png

\paragraph{Formalization target.}
create a compiling Lean file with sorry bodies at `PhyXMiniProblems/problem\_phyx\_mini\_0826.lean`. The Lean declarations must preserve the physical quantities, dimensions or dimensional roles, figure labels, governing-law hypotheses, and final relation expressed by this problem.
Use Mathlib/Physlib names found through LeanExplore where available. If a physics API is missing, introduce faithful local abstractions rather than scalar placeholder aliases.

\begin{theorem}[Physics formalization target]
\label{thm:physics:phyx_mini_0826:target}
\lean{PhyXMiniProblems.ProblemPhyXMini0826.finalCommonAngularVelocity_formula_and_choice_B}
\uses{def:physics:phyx-mini-0826:phyxminiproblems-problemphyxmini0826-angularvelocityinradianspersecond, def:physics:phyx-mini-0826:phyxminiproblems-problemphyxmini0826-momentofinertiainkilogrammeterssquared, def:physics:phyx-mini-0826:phyxminiproblems-problemphyxmini0826-diskclutchsetup, def:physics:phyx-mini-0826:phyxminiproblems-problemphyxmini0826-matchesprimarydiskclutchfigure, def:physics:phyx-mini-0826:phyxminiproblems-problemphyxmini0826-matchesdiskclutchproblemdata, def:physics:phyx-mini-0826:phyxminiproblems-problemphyxmini0826-hasphysicaldiskclutchparameters, def:physics:phyx-mini-0826:phyxminiproblems-problemphyxmini0826-satisfiestorquefreeclutchangularmomentumlaws, def:physics:phyx-mini-0826:phyxminiproblems-problemphyxmini0826-recordeddatasetanswer, def:physics:phyx-mini-0826:phyxminiproblems-problemphyxmini0826-matchesanswerchoice}
The assigned autoformalize agent should translate this physics problem into Lean declarations in the covered file, with theorem and lemma proof bodies written as `by sorry`.
\end{theorem}
\begin{proof}
This is an autoformalization task, not a proof task. Produce statements that can later be proved by the physics prover without weakening the physical model.
\end{proof}

% --- Archon declaration topology begin ---
\section{Declaration topology}
\begin{definition}[moment Of Inertia Dimension]
  \label{def:physics:phyx-mini-0826:phyxminiproblems-problemphyxmini0826-momentofinertiadimension}
  \lean{PhyXMiniProblems.ProblemPhyXMini0826.momentOfInertiaDimension}
  The dimension of an axial moment of inertia, mass * length\textasciicircum{}2
\end{definition}
\begin{proof}
  This is the declared model definition of moment Of Inertia Dimension.
\end{proof}
\begin{definition}[angular Momentum Dimension]
  \label{def:physics:phyx-mini-0826:phyxminiproblems-problemphyxmini0826-angularmomentumdimension}
  \lean{PhyXMiniProblems.ProblemPhyXMini0826.angularMomentumDimension}
  \uses{def:physics:phyx-mini-0826:phyxminiproblems-problemphyxmini0826-momentofinertiadimension}
  The dimension of angular momentum, mass * length\textasciicircum{}2 / time
\end{definition}
\begin{proof}
  This is the declared model definition of angular Momentum Dimension.
\end{proof}
\begin{definition}[force Dimension]
  \label{def:physics:phyx-mini-0826:phyxminiproblems-problemphyxmini0826-forcedimension}
  \lean{PhyXMiniProblems.ProblemPhyXMini0826.forceDimension}
  The dimension of force, mass * length / time\textasciicircum{}2
\end{definition}
\begin{proof}
  This is the declared model definition of force Dimension.
\end{proof}
\begin{definition}[torque Dimension]
  \label{def:physics:phyx-mini-0826:phyxminiproblems-problemphyxmini0826-torquedimension}
  \lean{PhyXMiniProblems.ProblemPhyXMini0826.torqueDimension}
  \uses{def:physics:phyx-mini-0826:phyxminiproblems-problemphyxmini0826-forcedimension}
  The dimension of torque, mass * length\textasciicircum{}2 / time\textasciicircum{}2
\end{definition}
\begin{proof}
  This is the declared model definition of torque Dimension.
\end{proof}
\begin{definition}[Axial Angular Velocity]
  \label{def:physics:phyx-mini-0826:phyxminiproblems-problemphyxmini0826-axialangularvelocity}
  \lean{PhyXMiniProblems.ProblemPhyXMini0826.AxialAngularVelocity}
  A signed angular velocity component along the oriented common axis. Radians are dimensionless, so this has inverse-time dimension
\end{definition}
\begin{proof}
  This is the declared model definition of Axial Angular Velocity.
\end{proof}
\begin{definition}[Axial Moment Of Inertia]
  \label{def:physics:phyx-mini-0826:phyxminiproblems-problemphyxmini0826-axialmomentofinertia}
  \lean{PhyXMiniProblems.ProblemPhyXMini0826.AxialMomentOfInertia}
  \uses{def:physics:phyx-mini-0826:phyxminiproblems-problemphyxmini0826-momentofinertiadimension}
  A nonnegative scalar moment of inertia about the common axis
\end{definition}
\begin{proof}
  This is the declared model definition of Axial Moment Of Inertia.
\end{proof}
\begin{definition}[Axial Angular Momentum]
  \label{def:physics:phyx-mini-0826:phyxminiproblems-problemphyxmini0826-axialangularmomentum}
  \lean{PhyXMiniProblems.ProblemPhyXMini0826.AxialAngularMomentum}
  \uses{def:physics:phyx-mini-0826:phyxminiproblems-problemphyxmini0826-angularmomentumdimension}
  A signed angular-momentum component along the common axis
\end{definition}
\begin{proof}
  This is the declared model definition of Axial Angular Momentum.
\end{proof}
\begin{definition}[Axial Force]
  \label{def:physics:phyx-mini-0826:phyxminiproblems-problemphyxmini0826-axialforce}
  \lean{PhyXMiniProblems.ProblemPhyXMini0826.AxialForce}
  \uses{def:physics:phyx-mini-0826:phyxminiproblems-problemphyxmini0826-forcedimension}
  A signed force component along the common axis
\end{definition}
\begin{proof}
  This is the declared model definition of Axial Force.
\end{proof}
\begin{definition}[Axial Torque]
  \label{def:physics:phyx-mini-0826:phyxminiproblems-problemphyxmini0826-axialtorque}
  \lean{PhyXMiniProblems.ProblemPhyXMini0826.AxialTorque}
  \uses{def:physics:phyx-mini-0826:phyxminiproblems-problemphyxmini0826-torquedimension}
  A signed torque component about the common axis
\end{definition}
\begin{proof}
  This is the declared model definition of Axial Torque.
\end{proof}
\begin{definition}[nonnegative SIReadout]
  \label{def:physics:phyx-mini-0826:phyxminiproblems-problemphyxmini0826-nonnegativesireadout}
  \lean{PhyXMiniProblems.ProblemPhyXMini0826.nonnegativeSIReadout}
  Read a nonnegative dimensionful quantity in coherent SI units
\end{definition}
\begin{proof}
  This is the declared model definition of nonnegative SIReadout.
\end{proof}
\begin{definition}[signed SIReadout]
  \label{def:physics:phyx-mini-0826:phyxminiproblems-problemphyxmini0826-signedsireadout}
  \lean{PhyXMiniProblems.ProblemPhyXMini0826.signedSIReadout}
  Read a signed dimensionful quantity in coherent SI units
\end{definition}
\begin{proof}
  This is the declared model definition of signed SIReadout.
\end{proof}
\begin{definition}[angular Velocity In Radians Per Second]
  \label{def:physics:phyx-mini-0826:phyxminiproblems-problemphyxmini0826-angularvelocityinradianspersecond}
  \lean{PhyXMiniProblems.ProblemPhyXMini0826.angularVelocityInRadiansPerSecond}
  \uses{def:physics:phyx-mini-0826:phyxminiproblems-problemphyxmini0826-axialangularvelocity, def:physics:phyx-mini-0826:phyxminiproblems-problemphyxmini0826-signedsireadout}
  Radian-per-second readout of an axial angular velocity
\end{definition}
\begin{proof}
  This is the declared model definition of angular Velocity In Radians Per Second.
\end{proof}
\begin{definition}[moment Of Inertia In Kilogram Meters Squared]
  \label{def:physics:phyx-mini-0826:phyxminiproblems-problemphyxmini0826-momentofinertiainkilogrammeterssquared}
  \lean{PhyXMiniProblems.ProblemPhyXMini0826.momentOfInertiaInKilogramMetersSquared}
  \uses{def:physics:phyx-mini-0826:phyxminiproblems-problemphyxmini0826-axialmomentofinertia, def:physics:phyx-mini-0826:phyxminiproblems-problemphyxmini0826-nonnegativesireadout}
  Kilogram-metre-squared readout of an axial moment of inertia
\end{definition}
\begin{proof}
  This is the declared model definition of moment Of Inertia In Kilogram Meters Squared.
\end{proof}
\begin{definition}[angular Momentum In Kilogram Meters Squared Per Second]
  \label{def:physics:phyx-mini-0826:phyxminiproblems-problemphyxmini0826-angularmomentuminkilogrammeterssquaredpersecond}
  \lean{PhyXMiniProblems.ProblemPhyXMini0826.angularMomentumInKilogramMetersSquaredPerSecond}
  \uses{def:physics:phyx-mini-0826:phyxminiproblems-problemphyxmini0826-axialangularmomentum, def:physics:phyx-mini-0826:phyxminiproblems-problemphyxmini0826-signedsireadout}
  Kilogram-metre-squared-per-second readout of axial angular momentum
\end{definition}
\begin{proof}
  This is the declared model definition of angular Momentum In Kilogram Meters Squared Per Second.
\end{proof}
\begin{definition}[axial Force In Newtons]
  \label{def:physics:phyx-mini-0826:phyxminiproblems-problemphyxmini0826-axialforceinnewtons}
  \lean{PhyXMiniProblems.ProblemPhyXMini0826.axialForceInNewtons}
  \uses{def:physics:phyx-mini-0826:phyxminiproblems-problemphyxmini0826-axialforce, def:physics:phyx-mini-0826:phyxminiproblems-problemphyxmini0826-signedsireadout}
  Newton readout of a signed axial force
\end{definition}
\begin{proof}
  This is the declared model definition of axial Force In Newtons.
\end{proof}
\begin{definition}[axial Torque In Newton Meters]
  \label{def:physics:phyx-mini-0826:phyxminiproblems-problemphyxmini0826-axialtorqueinnewtonmeters}
  \lean{PhyXMiniProblems.ProblemPhyXMini0826.axialTorqueInNewtonMeters}
  \uses{def:physics:phyx-mini-0826:phyxminiproblems-problemphyxmini0826-axialtorque, def:physics:phyx-mini-0826:phyxminiproblems-problemphyxmini0826-signedsireadout}
  Newton-metre readout of a signed axial torque
\end{definition}
\begin{proof}
  This is the declared model definition of axial Torque In Newton Meters.
\end{proof}
\begin{definition}[Rotating Body]
  \label{def:physics:phyx-mini-0826:phyxminiproblems-problemphyxmini0826-rotatingbody}
  \lean{PhyXMiniProblems.ProblemPhyXMini0826.RotatingBody}
  The two coaxial rotating bodies named in the problem
\end{definition}
\begin{proof}
  This is the declared model definition of Rotating Body.
\end{proof}
\begin{definition}[Figure Stage]
  \label{def:physics:phyx-mini-0826:phyxminiproblems-problemphyxmini0826-figurestage}
  \lean{PhyXMiniProblems.ProblemPhyXMini0826.FigureStage}
  The two stages printed at the left of the supplied image
\end{definition}
\begin{proof}
  This is the declared model definition of Figure Stage.
\end{proof}
\begin{definition}[Screen Rotation Sense]
  \label{def:physics:phyx-mini-0826:phyxminiproblems-problemphyxmini0826-screenrotationsense}
  \lean{PhyXMiniProblems.ProblemPhyXMini0826.ScreenRotationSense}
  Rotation sense as it appears on the two-dimensional image
\end{definition}
\begin{proof}
  This is the declared model definition of Screen Rotation Sense.
\end{proof}
\begin{definition}[Pressing Force Label]
  \label{def:physics:phyx-mini-0826:phyxminiproblems-problemphyxmini0826-pressingforcelabel}
  \lean{PhyXMiniProblems.ProblemPhyXMini0826.PressingForceLabel}
  The two force symbols printed beside the bodies
\end{definition}
\begin{proof}
  This is the declared model definition of Pressing Force Label.
\end{proof}
\begin{definition}[Disk Clutch Figure Label]
  \label{def:physics:phyx-mini-0826:phyxminiproblems-problemphyxmini0826-diskclutchfigurelabel}
  \lean{PhyXMiniProblems.ProblemPhyXMini0826.DiskClutchFigureLabel}
  Text and mathematical labels visible in image 826.png
\end{definition}
\begin{proof}
  This is the declared model definition of Disk Clutch Figure Label.
\end{proof}
\begin{definition}[Disk Clutch Figure]
  \label{def:physics:phyx-mini-0826:phyxminiproblems-problemphyxmini0826-diskclutchfigure}
  \lean{PhyXMiniProblems.ProblemPhyXMini0826.DiskClutchFigure}
  \uses{def:physics:phyx-mini-0826:phyxminiproblems-problemphyxmini0826-rotatingbody, def:physics:phyx-mini-0826:phyxminiproblems-problemphyxmini0826-figurestage, def:physics:phyx-mini-0826:phyxminiproblems-problemphyxmini0826-screenrotationsense, def:physics:phyx-mini-0826:phyxminiproblems-problemphyxmini0826-pressingforcelabel, def:physics:phyx-mini-0826:phyxminiproblems-problemphyxmini0826-diskclutchfigurelabel}
  Qualitative evidence transcribed from the primary image. Screen rotation sense is kept distinct from signed physical axial direction because the two initial arrows are drawn on opposite exposed faces
\end{definition}
\begin{proof}
  This is the declared model definition of Disk Clutch Figure.
\end{proof}
\begin{definition}[Disk Clutch Setup]
  \label{def:physics:phyx-mini-0826:phyxminiproblems-problemphyxmini0826-diskclutchsetup}
  \lean{PhyXMiniProblems.ProblemPhyXMini0826.DiskClutchSetup}
  \uses{def:physics:phyx-mini-0826:phyxminiproblems-problemphyxmini0826-axialangularvelocity, def:physics:phyx-mini-0826:phyxminiproblems-problemphyxmini0826-axialmomentofinertia, def:physics:phyx-mini-0826:phyxminiproblems-problemphyxmini0826-axialangularmomentum, def:physics:phyx-mini-0826:phyxminiproblems-problemphyxmini0826-axialforce, def:physics:phyx-mini-0826:phyxminiproblems-problemphyxmini0826-axialtorque, def:physics:phyx-mini-0826:phyxminiproblems-problemphyxmini0826-rotatingbody, def:physics:phyx-mini-0826:phyxminiproblems-problemphyxmini0826-diskclutchfigure}
  Independent physical quantities for the initial bodies and coupled final assembly. In particular, finalCommonAngularVelocity is an unknown observable; it is not defined from conservation or from an answer choice
\end{definition}
\begin{proof}
  This is the declared model definition of Disk Clutch Setup.
\end{proof}
\begin{definition}[Matches Primary Disk Clutch Figure]
  \label{def:physics:phyx-mini-0826:phyxminiproblems-problemphyxmini0826-matchesprimarydiskclutchfigure}
  \lean{PhyXMiniProblems.ProblemPhyXMini0826.MatchesPrimaryDiskClutchFigure}
  \uses{def:physics:phyx-mini-0826:phyxminiproblems-problemphyxmini0826-diskclutchsetup}
  Facts read directly from the supplied BEFORE/AFTER image
\end{definition}
\begin{proof}
  This is the declared model definition of Matches Primary Disk Clutch Figure.
\end{proof}
\begin{definition}[Matches Disk Clutch Problem Data]
  \label{def:physics:phyx-mini-0826:phyxminiproblems-problemphyxmini0826-matchesdiskclutchproblemdata}
  \lean{PhyXMiniProblems.ProblemPhyXMini0826.MatchesDiskClutchProblemData}
  \uses{def:physics:phyx-mini-0826:phyxminiproblems-problemphyxmini0826-angularvelocityinradianspersecond, def:physics:phyx-mini-0826:phyxminiproblems-problemphyxmini0826-diskclutchsetup}
  Qualitative data stated in the problem. Positivity relative to the chosen axis encodes that both initial rotations have the same physical direction. The final angular velocity remains unconstrained here beyond being a common observable after frictional coupling
\end{definition}
\begin{proof}
  This is the declared model definition of Matches Disk Clutch Problem Data.
\end{proof}
\begin{definition}[Has Physical Disk Clutch Parameters]
  \label{def:physics:phyx-mini-0826:phyxminiproblems-problemphyxmini0826-hasphysicaldiskclutchparameters}
  \lean{PhyXMiniProblems.ProblemPhyXMini0826.HasPhysicalDiskClutchParameters}
  \uses{def:physics:phyx-mini-0826:phyxminiproblems-problemphyxmini0826-momentofinertiainkilogrammeterssquared, def:physics:phyx-mini-0826:phyxminiproblems-problemphyxmini0826-axialforceinnewtons, def:physics:phyx-mini-0826:phyxminiproblems-problemphyxmini0826-diskclutchsetup}
  Nondegeneracy conditions for the displayed physical branch. They ensure the sum I\_A + I\_B is nonzero but state no solved final angular velocity
\end{definition}
\begin{proof}
  This is the declared model definition of Has Physical Disk Clutch Parameters.
\end{proof}
\begin{definition}[Satisfies Torque Free Clutch Angular Momentum Laws]
  \label{def:physics:phyx-mini-0826:phyxminiproblems-problemphyxmini0826-satisfiestorquefreeclutchangularmomentumlaws}
  \lean{PhyXMiniProblems.ProblemPhyXMini0826.SatisfiesTorqueFreeClutchAngularMomentumLaws}
  \uses{def:physics:phyx-mini-0826:phyxminiproblems-problemphyxmini0826-angularvelocityinradianspersecond, def:physics:phyx-mini-0826:phyxminiproblems-problemphyxmini0826-momentofinertiainkilogrammeterssquared, def:physics:phyx-mini-0826:phyxminiproblems-problemphyxmini0826-angularmomentuminkilogrammeterssquaredpersecond, def:physics:phyx-mini-0826:phyxminiproblems-problemphyxmini0826-axialforceinnewtons, def:physics:phyx-mini-0826:phyxminiproblems-problemphyxmini0826-axialtorqueinnewtonmeters, def:physics:phyx-mini-0826:phyxminiproblems-problemphyxmini0826-diskclutchsetup}
  The axial pressing forces are equal and opposite, so each has zero moment about the common axis and the net external torque about that axis vanishes. Friction may torque each body, but it is internal to the two-body system. Consequently total axial angular momentum is conserved while the two moments of inertia add in the coupled assembly. These are governing relations. No field states the requested quotient for finalCommonAngularVelocity or identifies an answer choice
\end{definition}
\begin{proof}
  This is the declared model definition of Satisfies Torque Free Clutch Angular Momentum Laws.
\end{proof}
\begin{definition}[Answer Choice]
  \label{def:physics:phyx-mini-0826:phyxminiproblems-problemphyxmini0826-answerchoice}
  \lean{PhyXMiniProblems.ProblemPhyXMini0826.AnswerChoice}
  Labels of the four expressions printed with the problem
\end{definition}
\begin{proof}
  This is the declared model definition of Answer Choice.
\end{proof}
\begin{definition}[proposed Angular Velocity In Radians Per Second]
  \label{def:physics:phyx-mini-0826:phyxminiproblems-problemphyxmini0826-answerchoice-proposedangularvelocityinradianspersecond}
  \lean{PhyXMiniProblems.ProblemPhyXMini0826.AnswerChoice.proposedAngularVelocityInRadiansPerSecond}
  \uses{def:physics:phyx-mini-0826:phyxminiproblems-problemphyxmini0826-angularvelocityinradianspersecond, def:physics:phyx-mini-0826:phyxminiproblems-problemphyxmini0826-momentofinertiainkilogrammeterssquared, def:physics:phyx-mini-0826:phyxminiproblems-problemphyxmini0826-diskclutchsetup, def:physics:phyx-mini-0826:phyxminiproblems-problemphyxmini0826-answerchoice}
  The scalar angular-velocity expression printed beside each answer label. These definitions encode the supplied options; they do not define the independently stored final physical angular velocity
\end{definition}
\begin{proof}
  This is the declared model definition of proposed Angular Velocity In Radians Per Second.
\end{proof}
\begin{definition}[recorded Dataset Answer]
  \label{def:physics:phyx-mini-0826:phyxminiproblems-problemphyxmini0826-recordeddatasetanswer}
  \lean{PhyXMiniProblems.ProblemPhyXMini0826.recordedDatasetAnswer}
  \uses{def:physics:phyx-mini-0826:phyxminiproblems-problemphyxmini0826-answerchoice}
  The answer label recorded in the supplied dataset metadata
\end{definition}
\begin{proof}
  This is the declared model definition of recorded Dataset Answer.
\end{proof}
\begin{definition}[Matches Answer Choice]
  \label{def:physics:phyx-mini-0826:phyxminiproblems-problemphyxmini0826-matchesanswerchoice}
  \lean{PhyXMiniProblems.ProblemPhyXMini0826.MatchesAnswerChoice}
  \uses{def:physics:phyx-mini-0826:phyxminiproblems-problemphyxmini0826-angularvelocityinradianspersecond, def:physics:phyx-mini-0826:phyxminiproblems-problemphyxmini0826-diskclutchsetup, def:physics:phyx-mini-0826:phyxminiproblems-problemphyxmini0826-answerchoice}
  An unknown final readout agrees with the expression printed at a choice
\end{definition}
\begin{proof}
  This is the declared model definition of Matches Answer Choice.
\end{proof}
% --- Archon declaration topology end ---
% --- Archon physics formalization source end ---
